\documentclass[12pt, a4paper]{article}

% !TeX root = ../main.tex
\usepackage{a4wide}

\usepackage[utf8]{inputenc}

%\usepackage[ngerman]{babel}
\usepackage[english]{babel}

\usepackage[T1]{fontenc}
\usepackage{palatino}

\usepackage{graphicx}
\usepackage{caption}
\usepackage{url}
\usepackage{tocloft}
\usepackage{acronym}

\usepackage{mathpazo}
\usepackage{amsmath}
\usepackage{amsfonts}
\usepackage{adjustbox}

%\usepackage{subcaption}

\usepackage{hhline}
\usepackage{amssymb}
\usepackage{floatflt}
\usepackage{setspace}
\usepackage{float}
\usepackage{color}
\usepackage{listings}
\usepackage{array}
\usepackage{scrhack}
\usepackage{xcolor}
\usepackage{wrapfig}
\usepackage{hyperref}
\usepackage{url}
\usepackage{lmodern}
\usepackage{multirow}
\usepackage{subfig}
\usepackage{cleveref}
\usepackage{lipsum}


% !TeX root = ../main.tex
%%%%%%%%%%%%%%%%%%%%%%%%%%%%%%%%%%%%%%%%%%%%%%%%%%%%%%%%%%%%%%%%%%%%%%%%
% Data about you and the Document%
%%%%%%%%%%%%%%%%%%%%%%%%%%%%%%%%%%%%%%%%%%%%%%%%%%%%%%%%%%%%%%%%%%%%%%%%

% % Main Title of Document:
\newcommand{\myMaintitle}{Untersuchung und Aufbau einer DevOps Development Umgebung}

% % Sub Title of DocInput:
\newcommand{\mySubtitle}{Developing holistic software solutions through integration of existing individual solutions.}

% % Ihr Name:
\newcommand{\myName}{Henrik Gerdes}

% % Matrikelnummer:
\newcommand{\myMatrikel}{MatNr: 969272}

% % Ihr Geburtsort:
\newcommand{\brith}{Osnabrück}

% % Ihr Geburtsort:
\newcommand{\place}{Osnabrück}

% % Ihr Abgabedatum:
\newcommand{\submission}{\today}

% % Ihr Abgabedatum:
\newcommand{\mycourse}{Exposé für den B.Sc.}

% % Name des Betreuers/Erstprüfenden:
\newcommand{\fistSupervisor}{Dennis Ziegenhagen}
\newcommand{\secSupervisor}{Achim Hendrix}

% % In welchem Semester befinden Sie sich?
\newcommand{\mySemester}{6. Semester}

\title{\myMaintitle}

\author{\myName}
% !TeX root = ../main.tex
% % Zeilenabstand im Haupttext auf anderthalb-zeilig setzen
%\linespread{1.25}\selectfont

% Line spacing
%\onehalfspacing{}

%Pfad für Grafiken
\graphicspath{{fig/}}

%Styleregeln
\widowpenalty10000 % Vermeidet einzelne Zeilen eines Absatzes zu Beginn einer Seite
\clubpenalty10000 % Vermeidet einzelne Zeilen eines Absatzes am Ende einer Seite
\addtocontents{toc}{\protect\sloppy}
\setcounter{tocdepth}{3}


% % \sloppy bewirkt, dass Latex beim Blocksatz nicht über den rechten Rand hinausschreibt.
% % und dafür größere Lücken in einer Zeile in Kauf nimmt
\sloppy

% % Setzt Dokumenteigenschaften für PDFs, wenn das Paket 'hyperref' geladen wurde.
\hypersetup{pdftitle=\myMaintitle,pdfauthor=\myName,bookmarksopen=true}

%Source for picture captions
\newcommand{\source}[1]{\caption*{Source: {#1}} }

\newcommand{\code}[1]{\texttt{#1}}

\newcommand{\myparagraph}[1]{\paragraph{#1}\mbox{}\\}

\newcommand{\RM}[1]{\MakeUppercase{\romannumeral{} #1{}}}

\newcommand{\HRule}{\rule{\linewidth}{0.5mm}} % Defines a new command for horizontal


\definecolor{dkgreen}{rgb}{0,0.6,0}
\definecolor{gray}{rgb}{0.5,0.5,0.5}
\definecolor{mauve}{rgb}{0.58,0,0.82}

\lstset{ %
  language=Java,                  % the language of the code
  basicstyle=\footnotesize,       % the size of the fonts that are used for the code
  numbers=left,                   % where to put the line-numbers
  numberstyle=\tiny\color{gray},  % the style that is used for the line-numbers
  stepnumber=1,                   % the step between two line-numbers. If it's 1, each line
                                  % will be numbered
  numbersep=5pt,                  % how far the line-numbers are from the code
  backgroundcolor=\color{white},  % choose the background color. You must add \usepackage{color}
  showspaces=false,               % show spaces adding particular underscores
  showstringspaces=false,         % underline spaces within strings
  showtabs=false,                 % show tabs within strings adding particular underscores
  frame=single,                   % adds a frame around the code
  rulecolor=\color{black},        % if not set, the frame-color may be changed on line-breaks within not-black text (e.g. commens (green here))
  tabsize=4,                      % sets default tabsize to 2 spaces
  captionpos=b,                   % sets the caption-position to bottom
  breaklines=true,                % sets automatic line breaking
  breakatwhitespace=false,        % sets if automatic breaks should only happen at whitespace
  title=\lstname,                 % show the filename of files included with \lstinputlisting;
                                  % also try caption instead of title
  keywordstyle=\color{blue},          % keyword style
  commentstyle=\color{dkgreen},       % comment style
  stringstyle=\color{mauve}         % string literal style
}

%%%%%%%%%%%%%%%%%%%%%%%%%%%%%%%%%%%%%%%%%%%%%%%%%%%%%%%%%%%%%%%%%%%%%%%%%%%%%%%%%%%%%%%%%
%Examples
%%%%%%%%%%%%%%%%%%%%%%%%%%%%%%%%%%%%%%%%%%%%%%%%%%%%%%%%%%%%%%%%%%%%%%%%%%%%%%%%%%%%%%%%%
% \pdfmarkupcomment[markup=Squiggly,color=green]{with pdfcomment}{move to the front}.
% \pdfmarkupcomment[markup=StrikeOut,color=red]{stupid}{replace stupid with funny}
% \pdfmarkupcomment[markup=Highlight,color=yellow]{Of course, you can highlight complete sentences.}{Highlight}
% \pdfcomment[icon=Note,color=blue]{insert graphic!}

\begin{document}
\nocite{*}

\pagenumbering{gobble}
% !TeX root = ../main.tex
%%%%%%%%%%%%%%%%%%%%%%%%%%%%%%%%%%%%%%%%%
% Academic Title Page
% LaTeX Template
% Version 2.0 (17/7/17)
%
% This template was downloaded from:
% http://www.LaTeXTemplates.com
%
% Original author:
% WikiBooks (LaTeX - Title Creation) with modifications by:
% Vel (vel@latextemplates.com)
%
% License:
% CC BY-NC-SA 3.0 (http://creativecommons.org/licenses/by-nc-sa/3.0/)
%
% Instructions for using this template:
% This title page is capable of being compiled as is. This is not useful for
% including it in another document. To do this, you have two options:
%
% 1) Copy/paste everything between \begin{document} and \end{document}
% starting at \begin{titlepage} and paste this into another LaTeX file where you
% want your title page.
% OR
% 2) Remove everything outside the \begin{titlepage} and \end{titlepage}, rename
% this file and move it to the same directory as the LaTeX file you wish to add it to.
% Then add \input{./<new filename>.tex} to your LaTeX file where you want your
% title page.
%
%%%%%%%%%%%%%%%%%%%%%%%%%%%%%%%%%%%%%%%%%

%----------------------------------------------------------------------------------------
%	TITLE PAGE
%----------------------------------------------------------------------------------------
%Titelseite
\begin{titlepage}
	\centering
	\thispagestyle{empty}
	\begin{center}
	\includegraphics[width=0.9\textwidth]{uos.pdf}
	\end{center}
	\LARGE{\textsc{Institut für Informatik}}
	\vfill
	\textsc{\Large{\emph{\mycourse}}}\\[0.5cm]
	\HRule\\[0.4cm]
	\vspace{8mm}
	\huge{\textbf{{\fontfamily{ppl}\selectfont
	\myMaintitle}}}\\
	\HRule\\[0.4cm]
	\vspace{9mm}

	\begin{minipage}{0.4\textwidth}
		\begin{flushleft}
			\large
			\textit{Author}\\
			\textsc{\myName}\\ % Your name
			\textsc{\myMatrikel} % Your name
		\end{flushleft}
	\end{minipage}
	~
	\begin{minipage}{0.4\textwidth}
		\begin{flushright}
			\large
			\textit{Supervisor}\\
			\textsc{\fistSupervisor}\\ % Supervisor's name
			\textsc{\secSupervisor} % Supervisor's name
		\end{flushright}
	\end{minipage}

	\vspace{5cm}
	\large{\today}
	\vfill
	\end{titlepage}
	\newpage


\tableofcontents
\newpage
\newcounter{lastroman}
\setcounter{lastroman}{\value{page}}

\pagestyle{plain}
\pagenumbering{arabic}
% \maketitle

\section{Topic, Content and Structure}
\subsection{Topic}
With the continuing shift towards microservices in the software industry, more and more new tools are being developed to operate, scale and maintain services. The microservice culture stipulates new features are delivered to customers faster and more frequently, bringing developers and administrators closer together. The resulting principals for this short agile development cycle are combined under the concept of DevOps. While analyses, guidelines and experiences for the production operation of DevOps are available extensively, the benefits of their toolset for the use in development environments have hardly been analyzed yet.\newline
Individual Development environments are quiet unique for every developer, reflect personal preferences are often time-consuming and complex to setup. Microservices and DevOps tools can help to quickly set up a consistent development environments and yet provide the freedom to let developers use the preferred applications. Companies like Docker Inc. and GitHub/Microsoft are also pushing this concept forward with products like Development-Containers and CodeSpaces that allow users to quickly checkout code form any location even within a browser to test and apply changes to their applications.
\subsection{Goal}
The goal of the bachelor thesis is to provide an overview of development tools in a modern DevOps environment. There will be a description of the tools that are practically suitable to be used in a development environment alongside with the creation of a practical container-based developer environment example. The advantages and disadvantages of such an environment will be discussed and analyzed. If possible, a small evaluation based on practical experience will also be part of the theses knowledge gain.
\subsection{Research status}
\subsection{Concept}
The increasing demand and complexity of software alongside more frequent feature requests in the software industry resulted in creation and adoption of microservices. With the broader availability of containerization solutions and tools this trend continues. Containerization allows software to be independent of the host environment that runs the code. The container already ships with a pre- and well-defined runtime environment and therefore decouples the link between host and software. Multiple instances of the software can be deployed quickly and on demand without the need of tedious, time consuming and error-prone host configuration. However, these advantages also bring new challenges. Monolithic applications are developed and tested as one unit, components are tightly coupled so that errors between components can be found quickly. Microservices are often orchestrated and provide together one extensive application. Changes in one service may break other parts of the application, which often only becomes apparent during operation in the overall orchestrated system. Continuous Integration and Continuous Deployment are two concepts, which are increasingly adopted, that try to solve this problem. The creation and practice of these concepts are bundled under the principles of DevOps culture. While there are already comprehensive guides and materials for implementing CI/CD in DevOps, these mostly focus on the process of testing and delivery these resources fall short in the management of the development environment.\newline
Development environments are highly heterogeneous. Developers use different operating systems in various configurations, different versions of languages and their preferred tools. Further development environments have additional requirements like usage of demo data, source-maps, debug-symbols and opt-out security features. It is not uncommon that configuring such an environment takes a lot of time. The fact that these development endowments differ to the CI/CD builds may introduce unexpected behavior. This bachelor theses introduce the idea of using the same containerization approach and CI/CD pipeline to build well-defined development environments and evaluates the idea based on a practical example project.\newline
The process of building, application and usability of such a developer environment will be described in this bachelor thesis by means of an practical example.
\subsection{Motivation}
Development environments tend to be complex spicily for larger projects. Many projects require a specific version of a compiler, interpreter, library and other parts of the toolchain. The specifics on how to set up these environments differ based on the (flavor of the) operating system. However, when developing on different projects, different versions of compiler, interpreter and library are often required which can lead to conflicts. The existence of package manager tools like \ac*{python-venv} and \ac{nvm}, that try so solve the problem of incompatible libraries, poof that there is a need for dedicated clean project environments.\newline
Finally, setting up such an environment takes a lot of time and is a challenge for new team members. It may also act as a entry barrier in open source projects.\newline
Containerization is a approach to make the application independent of the host system and prevents configuration based errors. My hope is that the adoption of this concept in the development process reduces the setup time of an application as well as eliminate the before mentioned configuration errors. The approach of development containers makes it possible to quickly create a working development environment anywhere, even on a different host. The only thing the developer needs is an editor and a connection to that host. Services like codesandbox.io and GitHubs CodeSpaces are some of the first to implement these concepts for small projects.
\section{Outline draft}
This is a first draft of the theses outline.

\begin{enumerate}
    \item Introduction
    \begin{enumerate}
        \item Scope of the thesis
        \item Structural overview
        \item Goal of the theses
    \end{enumerate}
    \item Container and DevOps Basics
    \begin{enumerate}
        \item What is DevOps
        \item Container and Microservice Basics
        \item Containerization in development environments
    \end{enumerate}
    \item Concept of DevContainers
    \begin{enumerate}
        \item Pre requirements for DevContainers
        \item Description of a conceptual environment
        \item Available tools and resources
        \item Possible implementations approach's
        \item Strengths, weaknesses and limits
    \end{enumerate}
    \item Exemplary implementation process
    \begin{enumerate}
        \item Current state and goal
        \item Implementation approach
        \item The implementation process
        \item Encountered challenges and limits
        \item Final state
    \end{enumerate}
    \item Performance evaluation and analysis
    \begin{enumerate}
        \item Metrics and what to evaluate
        \item Evaluation and results
        \item Discussion of evaluation
    \end{enumerate}
    \item Future potential and outlook
    \item Conclusion
\end{enumerate}

\section{Schedule}
\begin{table}[H]
    \begin{tabular}{ll}
    \hline
    Period     & Task                                                                                                              \\ \hline\hline
    07.06.2021 & Further literature research                                                                                       \\ \hline
    18.06.2021 & Survey of the initial condition                                                                                   \\ \hline
    30.06.2021 & Implementation of the DevContainer setup                                                                          \\ \hline
    31.07.2021 & Extensive testing of the new environment and deployment at the company                                             \\ \hline
    15.08.2021 & Finalize the written part up to and including the exemplary implementation processes.                             \\ \hline
    31.08.2021 & Survey on practicability and satisfaction new development environment and writing of the evaluation \& conclusion \\ \hline
    17.09.2021 & Proofreading, printing and delivery                                                                               \\ \hline
    \end{tabular}
    \caption{Approximate schedule}
    \label{tab::time}
    \end{table}
\section{Literatur}
\bibliographystyle{IEEEtranSA}
\bibliography{bib/sources}
\newpage

% Anhang
\renewcommand{\thesubsection}{\Alph{subsection}}
\pagenumbering{Roman}
\setcounter{page}{\value{lastroman}}
\section*{Appendix}
\addcontentsline{toc}{section}{Appendix}
\listoffigures
%Abkürzungsverzeichnis
% !TeX root = ../main.tex
\newcommand{\abbr}{Abbreviations}
\subsection{Abbreviations}
%\addcontentsline{toc}{subsection}{Abbreviations}

\begin{acronym}[1234567890]		%[längste Abkürzung]
\setlength{\itemsep}{-\parsep}	% sorgt dafür, dass das Verzeichnis kompakt dargestellt wird.

\acro{BLEST}[BLEST]{BLock ESTimation}
\acro{CWND}[CWND]{Congestion Window}
\acro{DAPS}[DAPS]{Delay-Aware Packet Scheduler}
\acro{ECF}[ECF]{Earliest Completion First}
\acro{HTTP}[HTTP]{Hypertext Transfer Protocol}
\acro{ISP}[ISP]{Internet Service Provider}
\acro{IETF}[IETF]{Internet Engineering Task Force}
\acro{LRF}[LRF]{Lowest-RTT-First}
\acro{MSS}[MSS]{Maximum Segment Size}
\acro{MPTCP}[MPTCP]{Multipath TCP}
\acro{MPTCPSW}[MPTCP\textsubscript{SW}]{MPTCP's send window}
\acro{OTIAS}[OTIAS]{Out-of-Order Transmission for In-Order Arrival Scheduler}
\acro{RR}[RR]{Round Robin}
\acro{RTT}[RTT]{Round Trip Time}
\acro{SCTP}[SCTP]{Stream Control Transmission Protocol}
\acro{sRTT}[sRTT]{Smoothed RTT}
\acro{STTF}[STTF]{Shortest Transfer Time First}
\acro{TCP}[TCP]{Transmission Control Protocol}
\acro{VoIP}[VoIP]{Voice over IP}
\end{acronym}
% !TeX root = ../main.tex
\section*{Declaration of Authorship}

\vspace{5cm}

~\\
I hereby declare that the paper submitted is my own unaided work. I assure that I wrote this paper without using any other means and sources than those specified. As well as the sources used literally or analogously taken from the sources identified as such.

\vspace{3cm}
\begin{flushright}

\rule{8cm}{0.2mm} \\
Signature (\myName)
\end{flushright}

\vspace{2cm}
\place, the \submission{}

\end{document}